%% ============================================================
%%  RÉSUMÉ (Français)
%% ============================================================
\chapter*{Résumé}
\thispagestyle{plain}
\addcontentsline{toc}{chapter}{Résumé}

\noindent
Ce Projet de Master Professionnel présente la conception logique d'une couche
d'agents \gls{ia} destinée à la plateforme éducative SprintUp. Le travail livre
une architecture indépendante de toute pile technique particulière, que la
plateforme cible pourra ensuite implémenter sur sa propre infrastructure.

L'objet du document se concentre sur deux contributions de fond. La première
est un \textbf{mécanisme d'indexation} des activités existantes, qui permet à
un agent de vérifier qu'une activité similaire n'a pas déjà été produite avant
d'en générer une nouvelle. La seconde est la conception d'un \textbf{serveur
\gls{mcp}} qui expose les opérations de lecture et d'écriture du syllabus sous
la forme d'un catalogue d'outils typés, et la définition d'un \textbf{catalogue
d'agents} consommateurs de ce serveur, chacun restreint par une liste blanche
d'outils qui matérialise le principe du moindre privilège.

Le document s'organise en deux chapitres. Une introduction situe le projet
dans son contexte. Le chapitre~1 présente les motifs d'orchestration d'agents,
l'augmentation par recherche et le Model Context Protocol. Le chapitre~2
détaille la conception~: indexation vectorielle des activités, protocole
\gls{mcp} appliqué au système et catalogue des agents retenus. Une conclusion
générale clôt le document.

\vspace{0.5cm}
\noindent
\textbf{Mots-clés~:} IA générative, Model Context Protocol, indexation
sémantique, agents conversationnels, EdTech.
\newpage
